\section{Branching growth and the dilated horizon}\label{sec:branching}

Proposition~\ref{prop:nonzero-limit} gives a nonzero solution of the
linear equation \eqref{eq:limit-equation}. We prove that this solution
cannot persist for all positive dilated times under the cap
\(g\le(t+l)H\). The argument is the branching construction of
\cite[Section~9]{OpenAIPeriodic}, whose proof we give in full. Its inputs
are the limiting medium \eqref{eq:background}, the cap, and the local
transport formula established in Proposition~\ref{prop:blowup-limit}.

\begin{proposition}\label{prop:finite-dilated-horizon}
There are fixed choices
\[
K\ge1,\qquad 0<\eta\le\log2,\qquad
0<\alpha\le\tfrac14,\qquad 0<\gamma\le1
\]
with the following property. For every allowed geometric mesh with
lower endpoint \(l\in[0,1]\), the only nonnegative density solving
\eqref{eq:limit-equation} on \(0<t<\infty\), with locally integrable
gain and loss, the positive-time mild formula \eqref{dil:mild}, and the
bound \(g(t,x,v)\le(t+l)H(v)\), is zero.
Consequently every nonzero limit in
Propositions~\ref{prop:blowup-limit} and~\ref{prop:nonzero-limit} has
\(U<\infty\). The four parameters can be fixed before the thermal,
seed, bath, and final horizon parameters of the construction. Their
choice is independent of the mesh phase and of \(l\).
\end{proposition}

The proof compares growth with spatial spreading. In a constant medium,
we construct positive collision histories that multiply mass while
keeping all speeds bounded and prescribing a small mean transverse
velocity. Finitely many compact families of these histories remain
effective when the medium changes measurably within a small tolerance. Far enough from
the initial time, the actual background has precisely this form on
short intervals whose lengths are proportional to their starting times.
Iteration then makes mass grow faster than the cap allows in spatial
balls of radius proportional to time. All histories begin at a strictly
positive time, where the collision rates are bounded on bounded velocity
sets.

The decisive comparison is between a multiplication factor and the
volume of these balls. Their volume is of order \(t^3\), and the cap
contributes one further power of \(t\). We will obtain a factor greater
than \(3\) between times \(t\) and \((1+a_{\rm br})t\), where the fixed number
\(a_{\rm br}>0\) is small enough that \((1+a_{\rm br})^4<3\).

\subsection{Positive histories and their normalization}

For \(v,w\in\R^3\) and \(n\in S^2\), set
\begin{equation}\label{br:sphere-map}
\Phi(v,w,n)=\frac{v+w}{2}+\frac{|v-w|}{2}n,
\qquad \dd\varsigma(n)=\frac{\dd n}{4\pi}.
\end{equation}
Thus \(\varsigma\) is a probability measure. By
Lemma~\ref{lem:carleman}, a collision with incident velocities \(v,w\)
has rate \(2\pi|v-w|\), and the marginal law of either child is the
pushforward of \(\varsigma\) under \(\Phi(v,w,\cdot)\). The two marginal
laws agree because \(n\mapsto-n\) preserves \(\varsigma\). Consequently,
for a nonnegative velocity measure \(B\) and a density \(f\),
\begin{equation}\label{br:gain-mass}
\int_{\R^3}P(B,f)(v)\dd v=\int_{\R^3}\nu_B(v)f(v)\dd v.
\end{equation}
The coefficient \(2\) in the gain of \eqref{eq:limit-equation} therefore
counts both children, while the loss removes their parent once.

More generally, fix Borel representatives of the coefficients of a
two-beam medium
\[
\widetilde B(t,x)=\sum_{\sigma=\pm1}b_\sigma(t,x)
                  \delta_{w_\sigma(t,x)},\qquad
q_\sigma(t,x,v)=2\pi b_\sigma(t,x)|v-w_\sigma(t,x)|,
\qquad \nu=\sum_\sigma q_\sigma.
\]
All pointwise path conditions below refer to these representatives.
On \([a,b]\), a history with ordered collision times
\(a<\tau_1<\cdots<\tau_j<b\), signs \(\sigma_i\), and directions \(n_i\)
follows free flights between collisions and uses
\[
V_{\tau_i}=\Phi(V_{\tau_i-},w_{\sigma_i}(\tau_i,X_{\tau_i}),n_i).
\]
Its weight relative to \(\prod_i\dd\tau_i\dd\varsigma(n_i)\) is
\begin{equation}\label{br:history-weight}
\exp\left(-\int_a^b\nu(s,X_s,V_s)\dd s\right)
\prod_{i=1}^j2q_{\sigma_i}(\tau_i,X_{\tau_i},V_{\tau_i-}).
\end{equation}
For \(j=0\) the empty product is one.

\begin{lemma}\label{br:positive-histories}
Let \(0<a<b<\infty\). Suppose \(f\ge0\) solves
\(\D f=2P(\widetilde B,f)-\nu f\) in the local transport sense,
both collision terms are locally integrable, and the chosen rate
representatives \(q_\sigma\) are bounded on compact positive-time
phase-space sets. Start with any nonnegative bounded, integrable
portion of \(f(a)\). Integration of any measurable collection of bounded
finite histories with weights \eqref{br:history-weight} gives a measure
dominated by \(f(b)\).
\end{lemma}

\begin{proof}
Set \(G=2P(\widetilde B,f)\). In free coordinates
\(f^\sharp(t,x,v)=f(t,x+tv,v)\), the equation becomes
\[
\partial_t f^\sharp=(G-\nu f)^\sharp.
\]
On bounded tubes its right side is locally integrable, so \(f^\sharp\)
has absolutely continuous time sections for almost every \((x,v)\)
and local \(L^1\) time traces.
The rate bounds make \(\nu^\sharp\) integrable along every bounded
characteristic segment in a compact positive-time tube. The one-dimensional integrating-factor
formula therefore consists of the freely
transported initial density multiplied by its survival factor, plus the
time integral of the nonnegative gain with the same survival factor.
Repeated substitution into the gain gives \eqref{br:history-weight},
by \eqref{br:sphere-map} and \eqref{br:gain-mass}. Every remainder is
nonnegative. Restricting the histories, or the initial density, preserves
the lower bound. Tonelli's theorem justifies the substitutions. On bounded
portions with total gain rate at most \(C\), the \(j\)-collision term is
bounded by \(C^j(b-a)^j/j!\) times the initial mass.

These formulas apply to point-velocity partners. Indeed
Lemma~\ref{lem:carleman} gives a density for the gain against each such
partner, with mass \eqref{br:gain-mass}. Thus free flight and each gain
preserve absolute continuity when the starting measure has a density.
Null spatial sets in representatives of the coefficients do not affect
the integrated formula.
\end{proof}

For the background \eqref{eq:background}, \(|u_\sigma|\le1+\alpha\) and
\begin{equation}\label{br:global-rate}
0\le b_\sigma(t,x)\le\frac{\theta K}{t\alpha\gamma},
\qquad
\nu_B(t,x,v)\le\frac{C(1+|v|)}{t}\quad(t>0).
\end{equation}
At each mesh interface choose the value from the outer adjacent
annulus, and on the transverse axis set the medium to zero.
These are Borel representatives obeying the same bounds.
All bounded finite flights starting at
positive time have strictly positive survival. The cap and
Lemma~\ref{lem:sheet-bound} give locally integrable collision terms.
The preceding transport argument therefore supplies local \(L^1\)
time traces of \(g\) on bounded regions. For the extracted limits,
these agree with the representatives in
Proposition~\ref{prop:blowup-limit}. Testing
nonnegative compact functions and taking a time limit shows that
\(g(t)\le(t+l)H\) also holds at these traces. We use this representative.

\subsection{A reference flight with a positive potential}

We convert the branching history weights into expectations along a
single Markov flight with a positive potential. This is a jump-process
Feynman--Kac representation, related to the many-to-one formulas for
branching systems; see \cite[Section~4.1]{HarrisRoberts2017} for general
probabilistic context. The finite-history calculation below proves the
precise representation for the present collision law.

Fix
\[
D=5,\qquad E=\{v:|v|\le D\},\qquad
B_0=\delta_{e_3}+\delta_{-e_3},\qquad
\nu_0(v)=2\pi\bigl(|v-e_3|+|v+e_3|\bigr).
\]
Let \(J(v,\dd v')\) be the single-child law obtained by choosing a sign
with probabilities proportional to \(|v-\sigma e_3|\) and then using
\eqref{br:sphere-map}. Put \(p(v)=J(v,E)\).

\begin{lemma}\label{br:retention}
For every \(v\in E\), \(p(v)\ge25/26\).
The reference history sum restricted to velocities in \(E\) is represented
by a Markov flight with velocity jump rate
\(\lambda(v)=2p(v)\nu_0(v)\), jump law
\(J(v,\dd v')\ind_E(v')/p(v)\), and multiplicative weight
\[
\exp\left(\int_0^A \mathcal W(V_s)\dd s\right),
\qquad \mathcal W(v)=(2p(v)-1)\nu_0(v).
\]
One may take
\begin{equation}\label{br:potential-bounds}
0<c_0=\frac{48\pi}{13}\le \mathcal W(v),\qquad
\lambda(v)\le C_0=48\pi\quad(v\in E).
\end{equation}
\end{lemma}

\begin{proof}
Set \(s=|v|^2\), \(z=v_3\), and
\(R=\sqrt{(s+1)^2-4z^2}\). For either partner sign the squared child
speed has the law
\[
|v'|^2=\frac{s+1+RZ}{2},
\qquad Z\text{ uniform on }[-1,1].
\]
This follows by expanding the square in \eqref{br:sphere-map}: a
coordinate of a uniform point on \(S^2\) is uniform on \([-1,1]\).
The maximum is at most \(s+1\), so there is no exit for \(s\le24\).
For \(24<s\le25\), the exit probability, if positive, is
\[
\frac12+\frac{s-49}{2R}
\le\frac12+\frac{s-49}{2(s+1)}
=\frac{s-24}{s+1}\le\frac1{26}.
\]
The inequality uses \(R\le s+1\) and \(s-49<0\). If \(R=0\), the
squared speed is deterministic and at most 25. Mixing the signs
preserves the bound.

A Markov history of rate \(\lambda\) has the survival factor
\(e^{-\int\lambda}\), and one rate times jump-law factor per collision.
Multiplication by \(e^{\int(\lambda-\nu_0)}\) changes this exactly to
\eqref{br:history-weight}, restricted to \(E\). Both history sums
converge absolutely on finite intervals because their rates are bounded.
Finally \(4\pi\le\nu_0(v)\le24\pi\) on \(E\), and
\(2p(v)-1\ge12/13\), giving \eqref{br:potential-bounds}.
\end{proof}

\subsection{Three collisions give a volume minorization}

The positive potential gives growth along every retained reference
path. To control where that growth occurs, we also need cancellation in
the integrated velocity. Write \(P_t(v,\cdot)\) for the transition law
of the preceding velocity process. We prove that these laws approach a
symmetric stationary distribution by establishing a Doeblin-type
minorization: every starting velocity has a common positive portion of
its transition law \cite{Doeblin1940}. Three collisions supply this
common portion.

\begin{lemma}\label{br:minorization}
There are \(r_0>0\) and \(\varepsilon_0>0\) such that
\begin{equation}\label{br:doeblin}
P_1(v,A)\ge\varepsilon_0|A\cap B(0,r_0)|\quad(v\in E)
\end{equation}
for every Borel set \(A\subset E\). For each \(\zeta>0\), there is also
\(\varepsilon_\zeta>0\) such that
\[
P_1(v,B(e_1,\zeta)\cap\{|v'|<D\})\ge\varepsilon_\zeta\quad(v\in E).
\]
\end{lemma}

\begin{proof}
Given \(v\in E\), choose a sign with \(|v-\sigma e_3|\ge1\).
Use this partner for the first collision, and select an output within
a small fixed distance \(\delta\) of \(\sigma e_3\). Its probability
has a uniform positive lower bound: the incident distance lies in
\([1,6]\), and a cap of fixed Euclidean radius \(\delta\) at that
endpoint has area probability bounded below uniformly on the
corresponding diameter spheres. Such first outputs lie strictly in \(E\).

Use partner \(-\sigma e_3\) next. If the first output is exactly
\(\sigma e_3\), the second sphere is the unit sphere; restrict its
output \(u\) near \(e_1\). For the third collision use partner \(e_3\).
The last output is
\[
T(u,n)=\frac{u+e_3}{2}+\frac{|u-e_3|}{2}n.
\]
At \(u=e_1\), \(n_0=-(e_1+e_3)/\sqrt2\), this output is zero.
The two angular derivatives in \(n\) span \((e_1+e_3)^\perp\).
The vector \(e_3\) is tangent to the preceding unit sphere at \(e_1\),
and direct differentiation gives
\begin{equation}\label{br:rank-derivative}
\partial_{e_3}T(e_1,n_0)=\frac{e_1+3e_3}{4},\qquad
(e_1+e_3)\cdot\partial_{e_3}T(e_1,n_0)=1.
\end{equation}
Thus the map of the four angular coordinates in the last two collisions
to the last velocity has rank three.

Take the two coordinates of \(n\) and the coordinate of \(u\) with
derivative \(e_3\) as distinguished variables, and hold the fourth
coordinate in a short interval. The inverse function theorem, with
this fourth coordinate and the first output as parameters, gives
neighborhoods on which the distinguished-variable map is one-to-one,
its Jacobian and inverse Jacobian are bounded, and its image contains
a fixed ball \(B(0,r_0)\). Shrink these neighborhoods and \(\delta\)
so this holds for every first output within \(\delta\) of
\(\sigma e_3\), for both signs. Sphere area in these coordinates has
a positive lower density. Change of variables in the three distinguished
coordinates, and integration in the fourth, give a fixed positive
multiple of volume on \(B(0,r_0)\).

Force the three times into three fixed, disjoint intervals inside
\((0,1)\), and prohibit other jumps. The first selected relative speed
is at least 1, the second is close to 2, and the third close to
\(\sqrt2\). In the reference Markov process the selected transition
intensity is
\[
4\pi|v-\sigma e_3|\,
\ind_E(\Phi(v,\sigma e_3,n))\dd\varsigma(n).
\]
It has positive uniform lower density on the patches just specified.
The no-additional-jump factor is at least \(e^{-C_0}\).
Integrating the times proves \eqref{br:doeblin}. Keeping only the first
two collisions, selecting a small patch near \(e_1\) in the second,
and prohibiting later jumps proves the last assertion.
\end{proof}

\begin{lemma}\label{br:mixing}
The reference process has a stationary probability law \(\pi\) with
\(\int v\dd\pi(v)=0\). There are \(C,c>0\) such that
\[
\sup_{v\in E}\|P_t(v,\cdot)-\pi\|_{\mathrm{TV}}\le Ce^{-ct},
\qquad
\sup_{v\in E}\mathbb E_v\left|
A^{-1}\int_0^A V_s\dd s\right|^2\le\frac C A\quad(A\ge1).
\]
\end{lemma}

\begin{proof}
The normalized measure on the right side of \eqref{br:doeblin} is a
common part of positive mass \(\epsilon\) in all rows of \(P_1\).
Removing this common part shows that the map \(\mathfrak m\mapsto\mathfrak m P_1\)
contracts total variation distance by \(1-\epsilon\). Probability
measures form a complete space in that metric, so iteration yields a
unique fixed point \(\pi\) and uniform geometric convergence.
Since \(P_t\) commutes with \(P_1\), \(\pi P_t\) is another fixed point
of \(P_1\), hence equals \(\pi\). Contractivity of Markov kernels gives
the estimate between integer times.

Negation of velocity interchanges the two signs and preserves \(E\),
so uniqueness of \(\pi\) makes it symmetric and its mean zero.
Consequently \(\sup_{v\in E}|\mathbb E_vV_t|\le Ce^{-ct}\).
The Markov property and bounded velocities imply
\[
|\mathbb E(V_s\cdot V_t)|\le Ce^{-c(t-s)}\quad(0\le s\le t)
\]
uniformly over starting states. Integrating this estimate on
\([0,A]^2\) bounds the second moment of \(\int_0^A V_s\dd s\) by \(CA\).
\end{proof}

\subsection{A displacement block and measurable perturbations}

Fix \(d=1/10\), and choose once and for all
\begin{equation}\label{br:kappa-choice}
0<\varkappa<\min\left\{\frac14,\frac{c_0}{2C_0},\frac{r_0}{4}\right\}.
\end{equation}
In particular a small ball about \(\varkappa e_1\) lies in the
minorization ball.

To repeat growth in the cold background, we will transport mass through
boxes whose transverse distance from the axis grows proportionally to time:
\[
\mathcal B_t=\{(x,v):|x-\varkappa te_1|\le d\varkappa t,\ |v|\le D\},
\qquad t>0.
\]
For a duration \(\Delta>0\), a displacement within
\(d\varkappa\Delta\) of \(\varkappa\Delta e_1\)
carries every starting point of \(\mathcal B_t\) into
\(\mathcal B_{t+\Delta}\), provided the final velocity remains in
\(E\). Mixing makes the integrated velocity small over most of a
reference block. We reserve a final fraction \(\varkappa\) of the block
for a flight with velocity close to \(e_1\), producing the displacement
\(\varkappa A e_1\). The positive potential will compensate for the cost
of prohibiting collisions during that final flight.

\begin{lemma}\label{br:reference-block}
There is a finite \(L\) such that, for all
\[
A\in I_L:=\left[\frac L{1+d},\frac L{1-d}\right],\qquad v\in E,
\]
reference branching histories with velocities in \(E\), final velocity
strictly inside \(E\), and displacement \(Z_A\) satisfying
\begin{equation}\label{br:strict-drift}
|Z_A-\varkappa Ae_1|<\frac34d\varkappa A
\end{equation}
have total weight greater than 8.
\end{lemma}

\begin{proof}
Put \(A_0=(1-\varkappa)A-1\). On the first interval of length \(A_0\),
require the reference Markov flight to have integrated velocity of
norm at most \(Ad\varkappa/4\). Lemma~\ref{br:mixing} and Chebyshev's
inequality give probability at least \(1/2\), uniformly in the start
when \(A\) is sufficiently large. Over the next unit interval require
the last velocity to lie in \(B(e_1,d/4)\). The conditional probability
is at least a fixed \(\varepsilon>0\), by
Lemma~\ref{br:minorization}. Require no jump during the final interval
of length \(\varkappa A\), whose conditional probability is at least
\(e^{-C_0\varkappa A}\).

The errors in displacement from \(\varkappa Ae_1\) on the three intervals
are at most \(Ad\varkappa/4\), \(D\), and \(Ad\varkappa/4\).
For \(A>4D/(d\varkappa)\) their sum is less than \(3Ad\varkappa/4\).
The final velocity lies strictly in \(E\).
Lemma~\ref{br:retention} bounds the branching weight below by
\[
\frac{\varepsilon}{2}\exp((c_0-C_0\varkappa)A).
\]
Its exponent is positive by \eqref{br:kappa-choice}. Choose \(L\)
large enough that this exceeds 8, and all the preceding lower bounds
on \(A\) hold throughout \(I_L\).
\end{proof}

The background coefficients jump across annular interfaces. To use the
reference growth there, we first restrict it to compact families with
bounded collision counts and strict velocity and displacement margins.
These margins allow comparison with a merely measurable medium: its
values along each actual path need only remain close to the constant
reference values.

\begin{lemma}\label{br:robust-block}
For these fixed \(\varkappa,L\), there are \(J_{\mathrm{br}}\in\N\) and
\(\varepsilon_*\in(0,1/4)\) with the following property. Suppose a
measurable medium obeys, throughout every speed-\(D\) excursion from
a specified starting point over a duration \(A\in I_L\),
\begin{equation}\label{br:medium-closeness}
|b_\sigma-1|\le\varepsilon_*,
\qquad |w_\sigma-\sigma e_3|\le\varepsilon_*
\quad(\sigma=\pm1).
\end{equation}
From each starting velocity in \(E\), histories with between \(1\) and
\(J_{\mathrm{br}}\) collisions, speed at most \(D\) on every free flight, every
postcollision velocity and the final velocity strictly inside \(E\), and
\begin{equation}\label{br:robust-drift}
|Z_A-\varkappa Ae_1|<d\varkappa A
\end{equation}
have total weight greater than 3. The initial free flight may have
speed \(D\). No continuity of the medium is needed.
\end{lemma}

\begin{proof}
Fix a reference start and duration \((v,A)\in E\times I_L\).
The zero-collision term has weight \(e^{-\nu_0(v)A}\le1\).
Discarding it from the family in Lemma~\ref{br:reference-block}
leaves weight greater than 7. Monotone convergence in the positive
history count permits restricting to finitely many counts
\(j\ge1\) while retaining weight greater than 7. Every postcollision
velocity lies strictly in \(E\) almost surely: the intersection of a
nondegenerate diameter sphere with \(\partial E\) has sphere area zero,
and a degenerate collision has zero rate. This includes the final
velocity. No inward margin is required for the initial velocity.
The displacement inequality
is strict. Ordered times may be kept apart and away from the endpoints.
For each retained collision count and sign sequence, parameterize the
ordered times as fractions of \(A\) and partition the angular parameter
space into finitely many measurable pieces lying in sphere coordinate
charts. Inner regularity of the
resulting finite parameter measures allows a further restriction to
compact sets with total weight greater than 6. These compact sets have
positive distances from the forbidden time coincidences and from every
postcollision speed boundary. They also have a positive margin in
\(|Z_A-\varkappa Ae_1|<d\varkappa A\), since the reference histories obey the
stronger inequality \eqref{br:strict-drift}. We call these finitely many
compact parameter sets a test for \((v,A)\).

The reference paths and weights are continuous in these finite
parameters and in \((v,A)\). The fixed compact tests therefore retain
weight greater than 5 and at least half their original positive
margins in a sufficiently small relative neighborhood of that start
and duration in \(E\times I_L\).
These neighborhoods cover the compact set
\(E\times I_L\). A finite subcover gives finitely many tests, a maximum
count \(J_{\mathrm{br}}\), and uniform positive margins.

Use the same initial velocity, signs, directions, and fractional times
for each perturbed history and its reference history. Their first
flights have speed at most \(D\), including boundary starts.
The sphere map satisfies
\[
|\Phi(v,w,n)-\Phi(\widetilde v,\widetilde w,n)|
\le |v-\widetilde v|+|w-\widetilde w|.
\]
Induction bounds the velocity error after \(j\) collisions by
\(j\varepsilon_*\) and the displacement error by \(AJ_{\mathrm{br}}\varepsilon_*\).
This comparison does not compare the medium at two nearby positions:
its value at the actual perturbed position is close to the same
constant reference partner. A sufficiently small \(\varepsilon_*\)
preserves all the strict velocity and displacement margins. The
actual paths thus have speed at most \(D\), so they stay within the
region on which closeness was assumed. Formally this follows by
induction up to the first possible failure; the preserved margins
exclude such a failure.

The relative-speed errors and full loss-rate errors on these histories
are bounded by \(C_{J_{\mathrm{br}}}\varepsilon_*\). The density error gives the same
bound for the gain-rate factors. The interval has bounded length,
so the error in the integrated loss and its exponential is bounded
by \(C_{J_{\mathrm{br}},L}\varepsilon_*\). Finite products of the bounded gain factors
then show that \eqref{br:history-weight} changes by at most
\(C_{J_{\mathrm{br}},L}\varepsilon_*\) uniformly on each test. The tests have finite
parameter measure. Decreasing \(\varepsilon_*\) makes their integrals
greater than 3. All coefficient estimates were pointwise; measurability
is sufficient.
\end{proof}

\subsection{Uniform choices of the construction parameters}

Lemma~\ref{br:robust-block} reduces the growth argument to a comparison
of the background with two constant beams. We now choose the
construction parameters to make that comparison uniform on the moving
boxes. Before doing so, we reserve a velocity tolerance that will also
allow positive mass to enter such a box. For this purpose we will steer
its velocity close to \(\varkappa e_1\), then let it fly freely until it
enters the moving box. There is
\(\alpha_{\mathrm{acc}}>0\), depending only on \(\varkappa\) and
Lemma~\ref{br:minorization}, with this property: after a first collision
puts a velocity sufficiently close to \(\sigma e_3\), two more
collisions, with partners within \(\alpha_{\mathrm{acc}}\) of
\(-\sigma e_3,e_3\), have positive angular probability of ending in
\begin{equation}\label{br:access-velocity}
B(\varkappa e_1,d\varkappa/4).
\end{equation}
Indeed the inverse-coordinate construction in that lemma, with first
output exactly \(\sigma e_3\), maps a positive-measure angular set into
\(B(\varkappa e_1,d\varkappa/8)\), which is inside \(B(0,r_0)\).
Keep a compact subset of positive measure in this set. Its outputs
have a positive distance from the complement of
\eqref{br:access-velocity}; the intermediate velocities and relative
speeds remain in the neighborhoods of that construction. Uniform
continuity of the two sphere maps preserves this for sufficiently
small perturbations of the first output and partners, including
measurable partners at the actual path positions. Take the smaller
tolerance for the two possible signs. We also require
\(\alpha_{\mathrm{acc}}\) to be less than half the permitted radius
of the first-output neighborhood.

The constants \(\varkappa,L,J_{\mathrm{br}},\varepsilon_*,\alpha_{\mathrm{acc}}\) are now
fixed and do not involve any construction parameter. Fix preliminary
bounds \(\eta\le1\), \(\alpha\le1\), \(K\ge1\). For \(t>0\), put
\begin{equation}\label{br:moving-box}
\Delta=\frac{L\varkappa t}{\theta K},
\qquad a_{\rm br}=\frac{L\varkappa}{\theta K}.
\end{equation}
At a starting point \(x_0=(y_0,z_0)\) in \(\mathcal B_t\),
\[
(1-d)\varkappa t\le|y_0|\le(1+d)\varkappa t.
\]
Every speed-\(D\) excursion over \([t,t+\Delta]\) satisfies
\[
\big||y|-|y_0|\big|\le D\Delta,\qquad
|z|\le d\varkappa t+D\Delta.
\]
Set \(\delta_K=DL/(\theta K(1-d))\). By increasing \(K\), impose
\(\delta_K\le1/4\) and \(a_{\rm br}\le1/4\). Then
\[
|y|\ge(1-\delta_K)(1-d)\varkappa t.
\]
For all sufficiently large \(t\), every such excursion is outside the
lower mesh hole, since \(l\le1\). On its active annulus
\(e^{-\eta}|y|\le\varrho\le|y|\), for every mesh phase. Thus, with a fixed
constant \(C_1\) depending only on \(d,D\),
\begin{equation}\label{br:comparison-quantities}
\begin{split}
\left|\frac{|y_0|}{\varrho}-1\right|&\le C_1(\eta+\delta_K),\\
\frac{\tau}{\varrho}&\le\frac{C_1}{\varkappa},\\
\frac{|z-\tau u_\sigma|}{\varrho}&\le\frac{C_1}{\varkappa}
\qquad(t\le\tau\le t+\Delta).
\end{split}
\end{equation}
The last estimate uses \(|u_\sigma|\le2\).

For a fixed starting point \(x_0\), set
\(\rho_0=\theta K/|y_0|\) and introduce the coordinates
\[
 s=\rho_0(\tau-t),\qquad \xi=\rho_0(x-x_0).
\]
Velocities are unchanged in these coordinates. The duration
\(\Delta\) is the same for every start in \(\mathcal B_t\), whereas
\(\rho_0\) varies with \(x_0\). This is the reason the reference block
was proved for an interval of durations: here its duration is exactly
\begin{equation}\label{br:duration-scaling}
A=\rho_0\Delta=\frac{L\varkappa t}{|y_0|}\in I_L.
\end{equation}
With \(S_\sigma=S(\gamma(z-\tau u_\sigma)/\varrho)\), the density ratio
in these units is
\begin{equation}\label{br:density-ratio}
\frac{b_\sigma(\tau,x)}{\rho_0}
=\frac{|y_0|}{\varrho}\,
\frac{S_\sigma}{1+(\tau\alpha\gamma/\varrho)S_\sigma}.
\end{equation}
Because \(1-\operatorname{sech}^2 h=\tanh^2h\le h^2\),
\eqref{br:comparison-quantities} gives, when the right-hand small
quantities are at most one,
\begin{equation}\label{br:uniform-medium-error}
\left|\frac{b_\sigma}{\rho_0}-1\right|
\le C_2\left(\eta+\delta_K+
              (\gamma/\varkappa)^2+\alpha\gamma/\varkappa\right),
\qquad |u_\sigma-\sigma|\le\alpha.
\end{equation}
Here \(C_2\) is independent of \(t,l\), and mesh phase.

Choose \(0<\eta\le\log2\) small, then
\[
0<\alpha<\min\{1/4,\varepsilon_*,\alpha_{\mathrm{acc}}\},
\]
then \(0<\gamma\le1\) small, so that the first, third, and fourth terms
on the right of \eqref{br:uniform-medium-error}, including their factor
\(C_2\), sum to less than \(\varepsilon_*/2\).
Finally increase \(K\) so the \(\delta_K\) term is less than
\(\varepsilon_*/2\), the preliminary restrictions hold, and
\begin{equation}\label{br:growth-exponent}
(1+a_{\rm br})^4<3.
\end{equation}
This establishes \eqref{br:medium-closeness} for all the scaled blocks,
with no choice depending on a limit, mesh, or solution.

These choices precede the remaining choices in the approximation.
Making \(\alpha\gamma\) small increases only fixed constants in
\eqref{br:global-rate} and the earlier comparison estimates.
Proposition~\ref{cold:control} is valid for every fixed
\(K,\eta,\alpha,\gamma\), after the thermal width is made sufficiently
small. Proposition~\ref{early:cap-selection} then permits the seed
exponent and the final short horizon to be chosen after those constants.
The later choices of bath, geometric scale, and horizon therefore
preserve the block estimates just proved.

\begin{lemma}[Growth in the limiting medium]\label{br:block-growth}
For the parameters just fixed, let \(a_{\rm br}=L\varkappa/(\theta K)\).
There is a time \(t_{\mathrm{br}}>0\), independent of the mesh and
\(l\in[0,1]\), such that every density satisfying the equation, cap,
and mild formula of Proposition~\ref{prop:finite-dilated-horizon} on
an interval containing \([t,(1+a_{\rm br})t]\), with \(t\ge t_{\mathrm{br}}\),
obeys, at its positive-time traces,
\begin{equation}\label{br:mass-growth}
 M((1+a_{\rm br})t)\ge3M(t),\qquad
 M(s):=\int_{\mathcal B_s}g(s,x,v)\dd x\dd v.
\end{equation}
\end{lemma}

\begin{proof}
Choose \(t_{\mathrm{br}}\) so that
\((1-\delta_K)(1-d)\varkappa t_{\mathrm{br}}>1\).
Then every speed-\(D\) excursion in
\eqref{br:moving-box} stays outside the hole for every
\(l\in[0,1]\) and every \(t\ge t_{\mathrm{br}}\).
Put \(\Delta=a_{\rm br}t\).
At each starting point of \(\mathcal B_t\), apply
Lemma~\ref{br:robust-block} using the simultaneous space--time scaling
\eqref{br:duration-scaling}. Translating back its displacement
restriction gives
\[
|x_{\mathrm{end}}-x_{\mathrm{start}}-\varkappa\Delta e_1|
<d\varkappa\Delta.
\]
The starting radius \(d\varkappa t\), followed by the triangle inequality,
places the endpoint in \(\mathcal B_{t+\Delta}\).
The velocity belongs to \(E\): the selected positive-count histories
have a strictly interior final velocity. Their first flight need
only have speed at most \(D\), exactly as in the medium comparison.
History weights are unchanged by this scaling: each gain rate is
divided by \(\rho_0\), each time differential is multiplied by
\(\rho_0\), and the integrated loss is invariant. The selected
histories therefore have total weight greater than 3 uniformly over
starts in \(\mathcal B_t\).

Use the kernel obtained by integrating all histories with
\(1\le j\le J_{\mathrm{br}}\) that satisfy the block's velocity and displacement
restrictions. The finite compact tests above certify its weight
greater than 3; no choice of a test for each starting point is needed.
For each integer \(m\ge1\), integrate this positive kernel against
\(\min\{g(t),m\}\ind_{\mathcal B_t}\). This is a bounded, integrable
portion of the trace, so Lemma~\ref{br:positive-histories} applies.
Letting \(m\to\infty\) and using monotone convergence gives
\eqref{br:mass-growth}.
\end{proof}

\subsection{Access to a late box and the growth contradiction}

We have obtained a uniform multiplication factor for mass already in a
moving box. The remaining point is to put some positive mass there.
The following argument needs no lower bound for the size of that mass:
any positive amount suffices for the eventual growth contradiction.

\begin{lemma}\label{br:accessibility}
For the fixed parameters above, let \(g\not\equiv0\) solve
\eqref{eq:limit-equation} on \(0<t<\infty\) with
\(0\le g\le(t+l)H\), for any allowed mesh and \(l\in[0,1]\).
Then \(\int_{\mathcal B_t}g(t,x,v)\dd x\dd v>0\) at arbitrarily large
finite times.
\end{lemma}

\begin{proof}
At some positive time \(a_0\), a bounded, integrable portion of
\(g(a_0)\) has positive mass. Such a portion is obtained by restricting
to bounded phase-space sets and truncating the density. The velocity axis has Lebesgue measure zero, so a
further portion of positive mass satisfies
\[
|x|\le R,\qquad |v|\le R,\qquad |\bar v|\ge\delta>0.
\]
For a sufficiently long but finite free flight of length \(s\),
\[
|y+s\bar v|\ge s\delta-R>l+2
\]
on this portion. Its survival is bounded below by a positive constant
by \eqref{br:global-rate}. Lemma~\ref{br:positive-histories} gives
positive mass in a bounded region strictly outside the hole at
time \(a_1=a_0+s\). This freely transported portion remains bounded
and integrable.

Choose \(h>0\) small so that paths whose initial speed is at most
\(R\) and whose later speeds are at most 3 remain in a bounded region
with \(|y|>l+1\) during \([a_1,a_1+h]\). Both sign densities have
strictly positive lower bounds there. To see this directly, the time,
axial coordinate, and transverse radius are bounded, the transverse
radius has a positive lower bound, and
\(e^{-\eta}|y|\le\varrho\le|y|\). The argument of \(S\) is therefore
bounded, its numerator is positive, and the denominator in
\eqref{eq:background} is bounded above. The rates on the selected
bounded-speed paths also have finite upper bounds.

Confine three collisions to this short interval. Choose a reference
sign with \(|v-\sigma e_3|\ge1\) for the first collision. The actual
relative speed is at least \(1-\alpha\ge3/4\).
Choose its output in a sufficiently small cap about the actual
partner. The cap probability has a positive lower bound for inputs
bounded by \(R\). Since the actual partner is within \(\alpha\) of
\(\sigma e_3\), this first output lies in the reserved first-output
neighborhood. Next use the signs \(-\sigma\) and \(+1\), and the
compact angular tests defining \(\alpha_{\mathrm{acc}}\). Their final
velocity lies in \eqref{br:access-velocity}. The second output stays
near \(e_1\), so the second and third relative speeds have positive
lower bounds; all postcollision speeds are at most 3.
Put the three times in disjoint subintervals and prohibit other
collisions. Positive lower bounds for the selected gain factors,
angular probabilities, time-interval lengths, and survival give
strictly positive total weight. The sphere-map comparison used here
requires only uniform partner closeness, and therefore remains valid
for measurable coefficients.

There is consequently a positive portion at time \(a_2=a_1+h\), with
bounded spatial support and \(|v-\varkappa e_1|<d\varkappa/4\).
Truncating its density gives a bounded, integrable subportion of
positive mass.
Let \(R_2\) bound \(|x-\varkappa a_2e_1|\) on that portion.
After free flight to any sufficiently large finite time \(t\),
\[
|x+(t-a_2)v-\varkappa te_1|
\le R_2+\frac{d\varkappa}{4}(t-a_2)<d\varkappa t.
\]
Its speed is strictly less than \(D\). For each finite \(t\), survival
remains strictly positive by \eqref{br:global-rate}. The zero-collision
term of Lemma~\ref{br:positive-histories} proves the assertion.
\end{proof}

\begin{proof}[Proof of Proposition~\ref{prop:finite-dilated-horizon}]
Suppose that a nonzero density \(g\) solves
\eqref{eq:limit-equation} for all positive times and obeys the cap.
Lemma~\ref{br:accessibility} gives a time \(t_0\ge t_{\mathrm{br}}\)
with positive mass \(m_0\) in \(\mathcal B_{t_0}\).

Apply Lemma~\ref{br:block-growth} successively.
At \(t_n=(1+a_{\rm br})^nt_0\), induction gives
\[
\int_{\mathcal B_{t_n}}g(t_n,x,v)\dd x\dd v\ge3^n m_0.
\]
The cap and \(l\le1\), on the other hand, give
\[
\int_{\mathcal B_{t_n}}g(t_n,x,v)\dd x\dd v
\le (t_n+1)\frac{4\pi}{3}(d\varkappa t_n)^3\int_E H(v)\dd v
\le C(1+t_n)^4.
\]
This contradicts \eqref{br:growth-exponent} as \(n\to\infty\).
The positive mass \(m_0\) and starting time may depend on the particular
density; their sizes are immaterial to this comparison. The parameters
governing the blocks were fixed before the mesh and \(l\) were chosen.
This rules out every nonzero global density in the proposition.
Applying this conclusion to the nonzero limits of
Proposition~\ref{prop:nonzero-limit} gives \(U<\infty\).
\end{proof}
